\documentclass[UTF8, fontset=macnew, a4paper, 12pt]{ctexart}

\usepackage{geometry}
\geometry{top=2.5cm, bottom=2.5cm, left=2.5cm, right=2.5cm}
\usepackage{amsmath, amssymb}
\usepackage{booktabs}
\usepackage{multirow}
\usepackage{tabularx}
\usepackage{graphicx}
\usepackage{caption}
\usepackage{xcolor}
\usepackage{float}
\usepackage{enumitem}
\usepackage{fancyhdr}
\usepackage{setspace}
\usepackage{pifont}
\usepackage[hidelinks]{hyperref}

\graphicspath{{../output/figures/}}
\captionsetup{font=small, labelfont=bf, labelsep=quad}
\setlist{itemsep=2pt, parsep=0pt, topsep=4pt}
\onehalfspacing

\ctexset{
  section/format = \raggedright\zihao{4}\bfseries,
  section/name   = {},
  section/number = \bignum{\value{section}}、,
  subsection/format = \raggedright\zihao{5}\bfseries,
  subsection/number = \chinese{subsection},
  subsection/aftername = 、,
  subsubsection/format = \raggedright\normalsize\bfseries,
}
% 節次採大寫數字（壹貳參…），與內文的「第壹節」及範例論文的體例一致；
% 小節仍用 \chinese（一二三…），對應內文的「之一」「之二」。
\newcommand{\bignum}[1]{\ifcase#1\or 壹\or 貳\or 參\or 肆\or 伍\or
  陸\or 柒\or 捌\or 玖\or 拾\fi}
\renewcommand{\thesection}{\bignum{\value{section}}}
% 交叉引用子節時補上節次，否則 \ref 只印出「三」而無從判斷屬於哪一節。
% \p@subsection 是 LaTeX 寫入 .aux 時加在 \thesubsection 前的前綴。
\renewcommand{\thesubsection}{\chinese{subsection}}
\makeatletter
\renewcommand{\p@subsection}{\thesection 節之}
\makeatother
\setcounter{secnumdepth}{2}   % subsubsection 不編號（標題已自帶（一）（二））

\pagestyle{fancy}
\fancyhf{}
\fancyhead[L]{\small 對數轉換偏誤的解析刻畫}
\fancyhead[R]{\small 理論推導（1006）}
\fancyfoot[C]{\small \thepage}
\renewcommand{\headrulewidth}{0.4pt}

% U+2460~2463（\cn{1}~\cn{4}）不在 xeCJK 認定的 CJK 範圍內，會落到拉丁字型而缺字，
% 故改用 pifont 的圈號。
\newcommand{\cn}[1]{\ding{\number\numexpr171+#1\relax}}
\newcommand{\edag}{e_0^{\dagger}}
\newcommand{\LCe}{\mathrm{LC}e^{\dagger}}
\newtheorem{prop}{命題}
\newtheorem{lem}{引理}
\newcommand{\Ex}{\mathbb{E}}

\begin{document}

\begin{center}
{\zihao{3}\bfseries 以解析偏誤校正修改小人口 Lee-Carter 模型的估計}\\[4pt]
{\zihao{5} Using Analytic Bias Correction to Modify the Estimation of
the Lee--Carter Model for Small Populations}\\[12pt]
{\zihao{5} （林子立）\quad 指導教授：余清祥\,教授}\\[4pt]
{\zihao{5} 國立政治大學統計學系}\\[4pt]
{\zihao{5} 中華民國 115 年 9 月 27 日}
\end{center}

\vspace{6pt}
\begin{center}
\begin{minipage}{0.92\textwidth}
\small
\textbf{摘要}\par\vspace{4pt}
Lee-Carter 模型的年齡水準參數 $\alpha_x$ 在小人口下有方向明確的偏誤，
既有研究以模擬量化此一現象並以修勻或參考人口加以緩解。本文改由解析的
途徑處理：在卜瓦松假設下，以 $c_0$ 替代零死亡格並取對數後，逐格偏誤有閉式
$b(\mu;c_0)=\Ex[\log\max(D,c_0)]-\log\mu$，而 $\hat\alpha_x$ 的偏誤恰為
該年齡各格 $b$ 的算術平均。此為恆等式而非漸近近似。該式的三項性質皆可
解析得出：$\mu\to0$ 時為 $\log(c_0/\mu)$、$\mu$ 大時為
$-1/(2\mu)-5/(12\mu^2)$、兩者之間有唯一變號點，$c_0=0.5$ 時
$\mu^\ast=0.9234$；且 $\mu^\ast$ 隨 $c_0$ 移動，亦即偏誤符號的分界由
分析者選定的替代值決定。以台灣全國女性 2001--2024 年資料的配適為真值
作模擬，該式預測逐年齡偏誤的相關係數為 $1.000$、$1.000$、$0.988$。
據此提出逐格扣除該偏誤的修改估計，使 $\alpha$ 偏誤中位數下降 $2.3$ 至
$7.2$ 倍，並在配合資訊加權後於三個人口規模的每一項指標上優於標準
Lee-Carter。本文同時指出該校正\textbf{完全不改善}時間參數的漂移項，
反過來證明漂移項的偏誤與對數轉換無關。
\par\vspace{6pt}
\textbf{關鍵詞}：Lee-Carter 模型、小人口、對數轉換偏誤、零格替代、
偏誤減少、最適加權
\end{minipage}
\end{center}

\vspace{10pt}

\section{緒論}

平均餘命是衡量地區發展與居民健康的重要指標，而其估計須以年齡別死亡率為
基礎。我國官方生命表只涵蓋全國與縣市層級，未及鄉鎮市區，原因不在資料
不全而在樣本規模：368 個鄉鎮市區中有 217 個的單一性別人數在兩萬以下。
在此規模下，死亡率較低的年齡組經常整年沒有死亡人數，
Lee and Carter（1992）模型（以下簡稱 LC 模型）的參數估計因而不穩。

人數少所帶來的困難通常被理解為估計值的震盪，但 Wang et al.（2018）與
Yue et al.（2019）指出，LC 模型的年齡參數在人數減少時還會出現
\textbf{有方向的偏誤}，且在人數二十萬以下時特別明顯。這項區別重要，
因為變異數大可藉重複抽樣察覺，偏誤卻在所有樣本上朝同一方向偏移，
單看估計結果無從發現，也無法靠加寬區間補救。既有的補救方向有二：
一是向外借資訊，如 Li and Lee（2005）以一組相似人口的共同趨勢為基準；
二是修勻，Yue et al.（2019）即以修勻修改 LC 的估計，在小人口上取得改善。

本文採第三個方向。既有研究對偏誤的刻畫是\textbf{模擬的}——在若干設定下
量出偏誤有多大——因而無法回答為何是這麼大，也無法在使用者拿到自己的
資料時給出一個數字。本文指出，在卜瓦松假設下這個偏誤有\textbf{閉式}，
且該閉式可用以修改估計。所處理的問題依序為三：偏誤能否寫成閉式、
閉式能否預測實測、以及能否據以建立扣除偏誤的估計量。

須先界定範圍。本文只處理年齡水準參數 $\alpha_x$，不處理年齡敏感度參數
$\beta_x$ 與時間參數 $\kappa_t$ 的漂移項。這不是取捨而是推導的結果之一：
第\ref{sec:modify}將顯示，一個確定性的期望死亡數函數可以移除 $\alpha_x$
的偏誤卻完全不動漂移項，此點反過來證明漂移項的偏誤不是對數轉換造成的。
換言之，本文的推導同時界定了自己的適用範圍，並指出小人口 LC 的偏誤
至少來自兩個獨立的機制。

\section{研究方法}\label{sec:method}

\subsection{模型與記號}

令 $x=1,\dots,A$ 為年齡組、$t=1,\dots,T$ 為年份。LC 模型設定
\begin{equation}\label{eq:lc}
\log m_{x,t}=\alpha_x+\beta_x\kappa_t,\qquad
\sum_x\beta_x=1,\quad \sum_t\kappa_t=0,
\end{equation}
其中 $m_{x,t}$ 為中央死亡率，識別條件用以固定尺度與位置。
觀測為死亡數 $D_{x,t}$ 與曝露數 $E_{x,t}$，本文設
\begin{equation}\label{eq:pois}
D_{x,t}\sim\text{Poisson}(\mu_{x,t}),\qquad \mu_{x,t}=E_{x,t}\,m_{x,t},
\end{equation}
各格獨立。式~\eqref{eq:pois} 是本文全部推導的基礎，其適當性於
第\ref{sec:disc}討論。

標準 LC 的作法是把死亡率觀測值取對數後作奇異值分解。但 $D_{x,t}=0$ 時
$\log(0/E)=-\infty$，故實務上必以某正值 $c_0$ 替代，慣例取 $c_0=0.5$。
令
\begin{equation}\label{eq:ell}
\ell_{x,t}=\log\frac{\max(D_{x,t},\,c_0)}{E_{x,t}},
\end{equation}
則標準 LC 的估計量為
\begin{equation}\label{eq:svd}
\hat\alpha_x=\frac{1}{T}\sum_{t}\ell_{x,t},\qquad
(\hat{\boldsymbol\beta},\hat{\boldsymbol\kappa})\ \text{取}\
\{\ell_{x,t}-\hat\alpha_x\}\ \text{的第一組奇異向量}.
\end{equation}
式~\eqref{eq:svd} 的第一式是本文的起點：$\hat\alpha_x$ 就是該列的算術
平均，在奇異值分解之前即已定出，因此它的期望值可以逐格計算。

\subsection{對數轉換偏誤的閉式}

考慮單一格，設 $D\sim\text{Poisson}(\mu)$、$\mu=Em$。該格觀測值的偏誤為
\[
\Ex\!\left[\log\frac{\max(D,c_0)}{E}\right]-\log\frac{\mu}{E}
=\Ex\big[\log\max(D,c_0)\big]-\log\mu,
\]
曝露數約掉，故偏誤只依賴 $\mu$ 與 $c_0$。按 $D$ 的取值展開期望值，
$D=0$ 時被替代為 $c_0$、$D=d\ge1$ 時取 $\log d$，得

\begin{prop}\label{prop:b}
設 $D\sim\text{Poisson}(\mu)$、$\mu>0$、$c_0>0$，則
\begin{equation}\label{eq:b}
b(\mu;c_0)\;\equiv\;\Ex\big[\log\max(D,c_0)\big]-\log\mu
\;=\;e^{-\mu}\log c_0+\sum_{d\ge1}\frac{e^{-\mu}\mu^{d}}{d!}\log d-\log\mu .
\end{equation}
級數絕對收斂，實作上取 $d\le\mu+12\sqrt{\mu}+12$ 即達機器精度。
\end{prop}

由式~\eqref{eq:svd} 與期望值的線性性質立得

\begin{prop}\label{prop:alpha}
在式~\eqref{eq:lc} 與~\eqref{eq:pois} 之下，
\begin{equation}\label{eq:alphabias}
\Ex[\hat\alpha_x]-\alpha_x=\frac{1}{T}\sum_{t=1}^{T} b(\mu_{x,t};c_0).
\end{equation}
\end{prop}

式~\eqref{eq:alphabias} 是\textbf{恆等式}，不是漸近近似或一階近似：
只要式~\eqref{eq:pois} 成立且 $\hat\alpha_x$ 依式~\eqref{eq:svd} 定義，
該式對任何 $T$ 與任何 $\mu_{x,t}$ 皆精確成立。這與小樣本偏誤文獻中常見的
$O(n^{-1})$ 型結果性質不同，原因是 $\hat\alpha_x$ 不含奇異值分解。
$\hat\beta_x$ 與 $\hat\kappa_t$ 則為奇異向量，其偏誤無法以同法寫出。

\subsection{閉式的三項性質}

\paragraph{（一）$\mu\to0$：零格替代造成的高估。}
注意 $\log 1=0$ 使 $d=1$ 項消失，故
\begin{equation}\label{eq:small}
b(\mu;c_0)=e^{-\mu}\Big[\log c_0+\tfrac{\mu^2}{2}\log 2+\cdots\Big]-\log\mu
=\log\frac{c_0}{\mu}-\mu\log c_0+O(\mu^2).
\end{equation}
$\mu\to0$ 時 $b\to+\infty$，速率 $\log(c_0/\mu)$。其意義直接：
$\mu<c_0$ 時該格幾乎必然為零，以 $c_0$ 替代即宣稱該格死亡率為 $c_0/E$，
而真值 $\mu/E<c_0/E$，故必然高估。$c_0=0.5$、$\mu=0.05$ 時
式~\eqref{eq:small} 給 $2.33724$，式~\eqref{eq:b} 的精確值 $2.337237$。

\paragraph{（二）$\mu$ 大：Jensen 不等式造成的低估。}
令 $u=(D-\mu)/\mu$，則 $\log D=\log\mu+\log(1+u)$。卜瓦松的中心動差給出
$\Ex[u]=0$、$\Ex[u^2]=1/\mu$、$\Ex[u^3]=1/\mu^2$、
$\Ex[u^4]=3/\mu^2+1/\mu^3$，代入
$\log(1+u)=u-\tfrac{u^2}{2}+\tfrac{u^3}{3}-\tfrac{u^4}{4}+\cdots$ 得
\begin{equation}\label{eq:jensen}
b(\mu;c_0)=-\frac{1}{2\mu}-\frac{5}{12\mu^{2}}+O(\mu^{-3}),
\end{equation}
其中 $-\tfrac{5}{12}=\tfrac13-\tfrac34$ 由第三、四項合併。$c_0$ 在此消失，
因為 $P(D=0)$ 可忽略。表~\ref{tab:asym} 顯示其精度。

\begin{table}[H]
\centering
\caption{式~\eqref{eq:jensen} 的精度（$c_0=0.5$）}
\small
\begin{tabular}{rrrrr}
\toprule
$\mu$ & 式~\eqref{eq:b} 精確值 & 一階 $-\frac{1}{2\mu}$ & 二階 式~\eqref{eq:jensen} & 二階誤差\\
\midrule
10   & $-0.0552395$ & $-0.0500000$ & $-0.0541667$ & $-1.1\times10^{-3}$\\
20   & $-0.0261518$ & $-0.0250000$ & $-0.0260417$ & $-1.1\times10^{-4}$\\
50   & $-0.0101730$ & $-0.0100000$ & $-0.0101667$ & $-6.4\times10^{-6}$\\
100  & $-0.0050424$ & $-0.0050000$ & $-0.0050417$ & $-7.7\times10^{-7}$\\
500  & $-0.0010017$ & $-0.0010000$ & $-0.0010017$ & $-6.0\times10^{-9}$\\
\bottomrule
\end{tabular}
\label{tab:asym}
\end{table}

\paragraph{（三）唯一變號點，且其位置由 $c_0$ 決定。}
由（一）與（二），$b$ 在 $\mu\to0^{+}$ 為 $+\infty$、在 $\mu\to\infty$ 時
由負側趨近 $0$。數值上 $b$ 先遞減、於 $\mu=2.2886$ 取極小值 $-0.19104$ 後
遞增趨近 $0^{-}$，故零點唯一。$c_0=0.5$ 時 $\mu^\ast=0.9234$。
亦即 $\alpha_x$ 偏誤的\textbf{符號由該年齡的單年期望死亡數決定}：
低於 $0.92$ 人者偏高、高於者偏低。

更值得注意的是 $\mu^\ast$ 對 $c_0$ 的依賴（表~\ref{tab:mustar}）。
偏誤符號的分界線\textbf{不是資料的性質，而是分析者選定替代值的後果}。
$c_0=0.5$ 這個慣例並無統計上的理由，而它決定了哪些年齡被高估。
此點在既有文獻未見討論，其後果有二：不同研究若採不同 $c_0$，
所報告的偏誤方向可能相反而不可直接比較；而 $c_0$ 雖可被選來使某個年齡
無偏，卻無法使所有年齡同時無偏——因 $\mu_{x,t}$ 跨年齡橫跨三個數量級。
本文因此不調整 $c_0$，而直接扣除 $b$。

\begin{table}[H]
\centering
\caption{變號點 $\mu^\ast$ 隨零格替代值 $c_0$ 的變化}
\small
\begin{tabular}{lrrrr}
\toprule
$c_0$ & 0.1 & 0.3 & 0.5 & 1.0\\
\midrule
$\mu^\ast$ & 0.134 & 0.533 & \textbf{0.923} & 1.509\\
\bottomrule
\end{tabular}
\label{tab:mustar}
\end{table}

\section{以台灣資料驗證閉式}\label{sec:valid}

式~\eqref{eq:alphabias} 雖為恆等式，但其推導依賴卜瓦松假設，
且實際估計流程含奇異值分解與識別條件的正規化，故須確認它在完整流程下成立。

驗證的設計如下。資料取自 Human Mortality Database 的台灣女性
2001--2024 年死亡數與曝露數，聚合為 22 個五齡組。以全國資料的卜瓦松配適
為真值 $(\alpha^{0},\beta^{0},\kappa^{0})$，按各年齡的曝露結構縮放至
人數 $N\in\{10^4,5\times10^4,2\times10^5\}$，依式~\eqref{eq:pois} 產生
死亡數，重複 100 次，固定亂數種子。以標準 LC 估計後，取各年齡偏誤的
中位數與式~\eqref{eq:alphabias} 的預測值比較。

\begin{table}[H]
\centering
\caption{閉式預測與模擬實測的對照（標準 LC，重複 100 次）}
\small
\begin{tabular}{lrrr}
\toprule
 & $N=10^4$ & $N=5\times10^4$ & $N=2\times10^5$\\
\midrule
逐年齡的相關係數 & \textbf{1.000} & \textbf{1.000} & \textbf{0.988}\\
迴歸斜率（實測對預測） & 1.001 & 0.990 & 0.962\\
殘差量級 & $\sim$0.01 & $\sim$0.01 & $\sim$0.02\\
\bottomrule
\end{tabular}
\label{tab:valid}
\end{table}

表~\ref{tab:valid} 顯示兩者幾乎重合，殘差落在蒙地卡羅誤差的量級
（重複 100 次時 $\alpha$ 偏誤的蒙地卡羅誤差約一成）。
表~\ref{tab:validage} 列出逐年齡的對照，並同時展示變號點：
期望死亡數 $0.9$ 人的 20 至 24 歲組偏誤已接近零，$1.2$ 人的零歲組
已轉為負，與 $\mu^\ast=0.9234$ 吻合。

\begin{table}[H]
\centering
\caption{逐年齡對照（$N=5\times10^4$，選列）}
\small
\begin{tabular}{lrrrr}
\toprule
年齡組 & 單年期望死亡數 & 閉式預測 & 模擬實測 & 差\\
\midrule
0      & 1.2  & $-0.0753$ & $-0.0692$ & $+0.0061$\\
1--4   & 0.3  & $0.7283$  & $0.7135$  & $-0.0148$\\
5--9   & 0.3  & $0.8689$  & $0.8658$  & $-0.0031$\\
10--14 & 0.3  & $0.8341$  & $0.8262$  & $-0.0079$\\
15--19 & 0.5  & $0.3298$  & $0.3386$  & $+0.0088$\\
20--24 & 0.9  & $0.0415$  & $0.0529$  & $+0.0114$\\
45--49 & 7.3  & $-0.0793$ & $-0.0831$ & $-0.0038$\\
70--74 & 48.8 & $-0.0108$ & $-0.0135$ & $-0.0027$\\
85--89 & 73.9 & $-0.0069$ & $-0.0061$ & $+0.0008$\\
\bottomrule
\end{tabular}
\label{tab:validage}
\end{table}

此一驗證把先前只能以模擬描述的現象變成一個恆等式：標準 LC 的 $\alpha_x$
偏誤\textbf{完全}由各格期望死亡數的一個確定性函數決定，不含未知成分。
既然如此，它就可以被算出來並扣掉。

\section{以解析偏誤校正修改 LC 的估計}\label{sec:modify}

\subsection{修改後的估計流程}

校正的對象是\textbf{格}而非參數，使校正同時流入三個參數。
困難在於 $b$ 依賴未知的 $\mu_{x,t}$，故須迭代：

\begin{enumerate}[label=\arabic*.]
\item 以標準 LC 起始，得 $(\hat{\boldsymbol\alpha},\hat{\boldsymbol\beta},\hat{\boldsymbol\kappa})$。
\item 計算 $\hat\mu_{x,t}=E_{x,t}\exp(\hat\alpha_x+\hat\beta_x\hat\kappa_t)$。
\item 以 $\ell_{x,t}-b(\hat\mu_{x,t};c_0)$ 取代 $\ell_{x,t}$，依式~\eqref{eq:svd} 重配適。
\item 回步驟 2，至參數變動小於容忍值。實測收斂約需 3 至 6 次。
\end{enumerate}

在偏誤減少的文獻中，Firth（1993）把方法分為兩類：\textbf{事後修正}
先算出最大概似估計值再扣除其偏誤的估計，此即 Cox and Snell（1968）；
\textbf{預先修正}則修改估計方程本身，此即 Firth 的 Jeffreys 懲罰。
本文的作法屬前者，但施加對象不同，而這個不同是決定性的。
Cox--Snell 修正最大概似估計量的 $O(n^{-1})$ 偏誤，其前提是最大概似
估計值\textbf{存在且有限}；而小人口 LC 的核心困難正是零格導致的完全分離
（Albert and Anderson 1984），卜瓦松最大概似的 $\hat\alpha_x$ 可為
$-\infty$（本文實測人數一萬時 5 至 9 歲組為 $-13.03$），
故 Cox--Snell 在此無從施加。式~\eqref{eq:b} 則完全不涉及最大概似估計量，
只用到 $D$ 的邊際分布，因此可在 Cox--Snell 失效處使用。
其代價是依賴卜瓦松假設，而 Firth 的懲罰只依賴概似的形式。

\subsection{加權：更換損失或校正之後該用什麼權重}

由診斷可知各年齡的噪音變異數差異達三個數量級，故加權為必要。
但「該用什麼權重」在更換損失函數後並非自明，此處一併推導。
設損失為 $\rho$、$\psi=\rho'$、逐格權重 $w$，估計方程為
$\sum_{x,t}w\,\psi(u)\,\partial u/\partial\theta=0$。
三明治變異數為 $A^{-1}BA^{-1}$，$A=\sum w\Ex[\psi'(u)]gg^{\top}$、
$B=\sum w^{2}\Ex[\psi(u)^{2}]gg^{\top}$，最小化之得
\begin{equation}\label{eq:wopt}
w\;\propto\;\Ex[\psi'(u)]\big/\Ex[\psi(u)^{2}] .
\end{equation}
設 $u\sim N(0,\sigma^2)$，三種損失的特例見表~\ref{tab:wcases}。
第一列即加權最小平方的標準結果；第二列則是 Koenker（2005，第 5.3 節）
定理 5.1 的內容——對檢查函數而言權重應取該分位處的\textbf{局部密度}，
而非標準差的倒數。

\begin{table}[H]
\centering
\caption{三種損失的變異數最適權重（$\beta_z=\Ex[\min(Z^2,z^2)]$，$Z\sim N(0,1)$）}
\small
\begin{tabular}{lllll}
\toprule
損失 & $\psi(u)$ & $\Ex[\psi'(u)]$ & $\Ex[\psi(u)^2]$ & 最適權重\\
\midrule
最小平方 & $u$ & 1 & $\sigma^2$ & $1/\sigma^2$\\
檢查函數 & $\tau-\mathbf{1}(u<0)$ & $f(0)$ & $\tau(1-\tau)$ & $f(0)$\\
Huber$(c)$ & $\min(|u|,c)\,\mathrm{sgn}(u)$ & $2\Phi(c/\sigma)-1$
  & $\sigma^2\beta_{c/\sigma}$ & $\dfrac{2\Phi(c/\sigma)-1}{\sigma^2\beta_{c/\sigma}}$\\
\bottomrule
\end{tabular}
\label{tab:wcases}
\end{table}

本問題的關鍵是 $\sigma_x$ 隨年齡變動。由卜瓦松的 delta 方法
$\mathrm{Var}(\log\hat m_{x,t})\approx1/\mu_{x,t}$，故
$\sigma_x=\mu_x^{-1/2}$、$c/\sigma_x=c\sqrt{\mu_x}$，代入得
\begin{equation}\label{eq:woptmu}
w_x\;\propto\;\mu_x\,\frac{2\Phi(c\sqrt{\mu_x})-1}{\beta_{c\sqrt{\mu_x}}} .
\end{equation}
兩個極限說明其結構：$c\sqrt{\mu}\to\infty$ 時 $w\propto\mu$（最小平方）；
$c\sqrt{\mu}\to0$ 時 $2\Phi(z)-1\approx z\sqrt{2/\pi}$、$\beta_z\approx z^2$，
故 $w\propto\sqrt{\mu}$（檢查函數，因 $\log\hat m$ 在中位數處的密度為
$\sqrt{\mu/2\pi}$）。故最適權重在 $\sqrt{\mu}$ 與 $\mu$ 之間內插，
而\textbf{內插的變數是 $c\sqrt{\mu}$ 而非 $c$}。
據此，固定冪次 $w=\mu^{p}$ 理論上不可能最適，因為同一筆資料的
$c\sqrt{\mu}$ 橫跨兩個數量級，式~\eqref{eq:woptmu} 的隱含冪次在
$c=1.345$ 時由 $0.61$ 漂移至 $0.87$。

\subsection{實測結果}

表~\ref{tab:estcmp} 以同一組亂數種子比較五個估計量。

\begin{table}[H]
\centering
\caption{解析校正與既有作法的比較（重複 100 次，取中位數）}
\small
\begin{tabular}{llrrrrr}
\toprule
$N$ & 估計量 & $\alpha$ 中位 & $\alpha$ 最大 & SSE($\beta$) & cor($\beta$) & 漂移偏誤\\
\midrule
\multirow{5}{*}{$10^4$}
 & 標準 LC          & 0.1612 & 2.3309 & 7.916 & $-0.074$ & 0.3436\\
 & 卜瓦松 MLE       & 0.0495 & 13.0317 & 15.525 & 0.118 & 0.1612\\
 & Firth            & \textbf{0.0248} & \textbf{0.4267} & 17.797 & 0.116 & 0.2237\\
 & 解析校正         & 0.0693 & 0.5008 & 18.048 & $-0.078$ & 0.3544\\
 & 解析校正 $+$ 加權 & 0.1749 & 0.7399 & \textbf{2.537} & 0.111 & \textbf{0.1884}\\
\midrule
\multirow{5}{*}{$5\times10^4$}
 & 標準 LC          & 0.0611 & 0.8658 & 4.449 & 0.037 & 0.3339\\
 & 卜瓦松 MLE       & 0.0073 & 0.2260 & 1.756 & 0.365 & 0.0492\\
 & Firth            & 0.0088 & \textbf{0.0680} & 1.410 & \textbf{0.397} & \textbf{0.0485}\\
 & 解析校正         & \textbf{0.0093} & 0.2434 & 8.834 & 0.198 & 0.3335\\
 & 解析校正 $+$ 加權 & 0.0110 & 0.2524 & \textbf{1.023} & 0.179 & 0.1381\\
\midrule
\multirow{5}{*}{$2\times10^5$}
 & 標準 LC          & 0.0231 & 0.1846 & 1.400 & 0.379 & 0.2117\\
 & 卜瓦松 MLE       & \textbf{0.0018} & 0.0361 & 0.346 & 0.596 & 0.0267\\
 & Firth            & 0.0025 & \textbf{0.0398} & \textbf{0.333} & \textbf{0.605} & \textbf{0.0262}\\
 & 解析校正         & 0.0032 & 0.1778 & 3.499 & 0.438 & 0.2747\\
 & 解析校正 $+$ 加權 & 0.0071 & 0.2897 & 0.832 & 0.213 & 0.1106\\
\bottomrule
\end{tabular}
\label{tab:estcmp}
\end{table}

成功之處有二。其一，$\alpha$ 偏誤中位數在三個規模下分別下降 $2.3$、
$6.6$、$7.2$ 倍；人數一萬時 $\alpha$ 最大偏誤由 $2.3309$ 降至 $0.5008$，
且決定性地優於卜瓦松最大概似的 $13.03$——後者因完全分離而崩壞，
這正是式~\eqref{eq:b} 不依賴最大概似這個性質的價值所在。
其二，加上資訊加權之後，該估計量在\textbf{三個規模的每一項指標上都優於
標準 LC}（$\beta$ 的平方誤差和 $7.92\to2.54$、$4.45\to1.02$、
$1.40\to0.83$；漂移項偏誤 $0.34\to0.19$、$0.33\to0.14$、$0.21\to0.11$），
且完全留在奇異值分解的架構內、沒有任何調校參數。

至於式~\eqref{eq:woptmu}，實測顯示\textbf{推導正確但無實務差異}。
以 $w\in\{1,\sqrt{\hat\mu},\hat\mu,\text{式~\eqref{eq:woptmu}}\}$ 比較，
三種加權版本彼此的差距皆在蒙地卡羅誤差之內——例如
$N=2\times10^5$、$c=1.345$ 時 $\beta$ 的平方誤差和分別為
$0.4051$、$0.4108$、$0.4017$——而「有加權對無加權」的差距則很大
（同一設定下無加權為 $1.0100$）。解釋是目標函數對權重的冪次相當平坦：
權重的作用是反映 $\sigma_x^2$ 跨三個數量級的落差，而 $\mu^{0.5}$、
$\mu^{1}$ 與式~\eqref{eq:woptmu} 都足以做到這件事，其間的差異屬二階。
實務上因此建議取最簡單的 $w=\hat\mu$。

另須指出式~\eqref{eq:wopt} 的一項界限。固定冪次的掃描顯示，
最適冪次只在常態近似成立處符合表~\ref{tab:wcases}：人數二十萬時
（各年齡 $\mu\ge1$）純檢查函數的最適冪次為 $0.50$、最小平方為 $1.00$，
完全一致；但人數一萬與五萬時無明顯型態。兩個原因：$\mu\ll1$ 時
$\ell_{x,t}$ 是有原子的離散分布，「局部密度」並不存在；
且式~\eqref{eq:wopt} 是\textbf{變異數}最適，而 $\alpha_x$ 在小人口下是
\textbf{偏誤主導}，變異數最適的權重不是均方誤差最適的權重。
這說明式~\eqref{eq:woptmu} 應與偏誤校正合用而非單獨使用。

\subsection{三項限制}

\paragraph{其一，純解析校正使 $\beta$ 變差。}
表~\ref{tab:estcmp} 中解析校正的 SSE($\beta$) 為 $18.05$、$8.83$、$3.50$，
均\textbf{大於}標準 LC 的 $7.92$、$4.45$、$1.40$。機制是逐格扣除
$b(\hat\mu_{x,t})$ 把 $\hat\mu$ 的估計誤差注入每一格，而 $\beta$ 的估計
依賴跨格的相對大小，因而對這種注入敏感，必須配合加權才能抵銷。
這是真實的代價，不是實作問題。

\paragraph{其二，漂移項不動。}
漂移項偏誤由 $0.3436$、$0.3339$、$0.2117$ 變為 $0.3544$、$0.3335$、
$0.2747$，即無改善。這是本文分工論最強的證據：既然一個確定性的 $\mu$
函數可以移除 $\alpha$ 的偏誤卻完全不動漂移項，那麼漂移項的偏誤就不可能是
對數轉換造成的。它來自等權重分解對異質變異的敏感，須以加權處理——
這也正是「解析校正 $+$ 加權」一列改善漂移項的原因。

\paragraph{其三，Firth 在 $\alpha$ 最大偏誤上仍勝出。}
Firth 的 $0.4267$、$0.0680$、$0.0398$ 全面優於解析校正的 $0.5008$、
$0.2434$、$0.1778$。理由與方法性質一致：事後扣除只能扣除\textbf{期望值}，
每一次實現的偏差仍在；Firth 在計數尺度上估計，根本不產生那筆錯誤資訊。
本文另掃描了部分扣除（係數 $0.8$、$0.6$），結果一律較完全扣除為差，
故\textbf{不存在過度校正}。

\section{討論與結論}\label{sec:disc}

本文把小人口 LC 模型 $\alpha_x$ 的偏誤由一個模擬觀察變成一個閉式。
三項結論如下。

第一，偏誤有解析形式且該形式是恆等式。式~\eqref{eq:alphabias} 顯示
$\hat\alpha_x$ 的偏誤為各格 $b(\mu_{x,t};c_0)$ 的平均，而
式~\eqref{eq:b} 給出 $b$ 的閉式。其三項性質——小 $\mu$ 的
$\log(c_0/\mu)$、大 $\mu$ 的 $-1/(2\mu)-5/(12\mu^2)$、
以及唯一變號點 $\mu^\ast$——把「哪些年齡被高估、哪些被低估、以及多少」
變成可計算的量。台灣資料的驗證給出 $1.000$、$1.000$、$0.988$ 的相關係數。
其中一項副產品值得向實務提醒：變號點 $\mu^\ast$ 隨零格替代值 $c_0$ 移動，
故偏誤的符號部分是分析者的選擇而非資料的性質。

第二，據此修改估計可觀改善 $\alpha_x$，但須配合加權。
修改後的估計使 $\alpha$ 偏誤中位數下降 $2.3$ 至 $7.2$ 倍，
配合資訊加權後在三個人口規模的每一項指標上優於標準 LC，
且無調校參數、不離開奇異值分解的架構。相對於 Yue et al.（2019）
以修勻修改估計，本文的修改不引入平滑參數，因而不需要選參；
但其代價是依賴卜瓦松假設。

第三，該校正完全不改善漂移項，此點本身即是結果。它證明小人口 LC 的偏誤
至少來自兩個獨立機制：$\alpha_x$ 源於對數轉換與零格替代，
可由本文的閉式處理；漂移項源於等權重分解對異質變異的敏感，
須以加權處理。兩者是獨立的病，也需要獨立的藥。

本文的推導留下四個缺口。第一，$\hat\beta_x$ 與 $\hat\kappa_t$ 沒有對應的
閉式，兩者為奇異向量，其偏誤須以矩陣擾動理論處理。
第二，式~\eqref{eq:b} 依賴卜瓦松假設；實際死亡數可能過度離散
（例如同一機構的死亡具相關性），此時式~\eqref{eq:jensen} 的一階項應為
$-\sigma_D^2/(2\mu^2)$ 而非 $-1/(2\mu)$，負二項假設下的對應閉式可同法
寫出但尚未做。第三，本文只處理 $c_0$ 為常數的情形，逐年齡的 $c_0$
是否有意義尚未探討。第四，本文以全國配適為真值作模擬，
尚未在真實的鄉鎮市區資料上驗證——而該驗證是本方法能否用於編製
鄉鎮市區生命表的最後一步。

\section*{參考文獻}
\addcontentsline{toc}{section}{參考文獻}
\begin{itemize}[leftmargin=2em, itemsep=2pt, parsep=0pt]
  \item[] Albert, A., and Anderson, J. A. 1984. ``On the Existence of Maximum
        Likelihood Estimates in Logistic Regression Models.''
        \textit{Biometrika} 71(1): 1-10.
  \item[] Cox, D. R., and Snell, E. J. 1968. ``A General Definition of Residuals.''
        \textit{Journal of the Royal Statistical Society, Series B} 30(2): 248-265.
  \item[] Firth, D. 1993. ``Bias Reduction of Maximum Likelihood Estimates.''
        \textit{Biometrika} 80(1): 27-38.
  \item[] Huber, P. J. 1964. ``Robust Estimation of a Location Parameter.''
        \textit{Annals of Mathematical Statistics} 35(1): 73-101.
  \item[] Koenker, R. 2005. \textit{Quantile Regression}. Cambridge: Cambridge
        University Press.
  \item[] Lee, R. D., and L. R. Carter. 1992. ``Modeling and Forecasting U.S.
        Mortality.'' \textit{Journal of the American Statistical Association}
        87(419): 659-671.
  \item[] Li, N., and R. Lee. 2005. ``Coherent Mortality Forecasts for a Group of
        Populations: An Extension of the Lee-Carter Method.''
        \textit{Demography} 42(3): 575-594.
  \item[] Wang, H. C., C. J. Yue, and C. T. Chong. 2018. ``Mortality Models and
        Longevity Risk for Small Populations.''
        \textit{Insurance: Mathematics and Economics} 78: 351-359.
  \item[] Yue, J. C., H. C. Wang, and T. Y. Wang. 2019. ``Using Graduation to
        Modify the Estimation of Lee-Carter Model for Small Populations.''
        \textit{North American Actuarial Journal}.
        doi:10.1080/10920277.2019.1650288
\end{itemize}

\end{document}
